% 出席確認クイズ 第12回　デジタル戦略とAI戦略
%   attendance/attendance12.html から生成（解説は本ガイド用に加筆）。

\quizq{1}{DX(デジタルトランスフォーメーション)の説明として適切なものはどれでしょうか?}%
  {ホームページを作ること}%
  {紙の書類をスキャンしてPDFにすること}%
  {パソコンを新しくすること}%
  {\correct{デジタル技術でビジネスモデルや組織を変革し、競争優位を築くこと}}%
  {正解は (d)。DX（デジタルトランスフォーメーション）は、デジタル技術を使ってビジネスモデルや組織のあり方を変革し、競争優位を築くことです。(a)(b)(c) は単なるデジタル化（デジタイゼーション）であり、DX の出発点にはなりますが DX そのものではありません。}

\quizq{2}{「プラットフォーム型ビジネス」の例として適切なものはどれでしょうか?}%
  {\correct{売り手と買い手を結びつけるオンラインマーケットプレイス}}%
  {単一店舗での対面販売}%
  {自社工場での製品製造}%
  {社内の定例会議}%
  {正解は (a)。プラットフォーム型ビジネスは、売り手と買い手のように複数の利用者グループを結びつけて価値を生みます。参加者が増えるほど価値が高まるネットワーク効果が特徴です。(b)(c)(d) は単独の企業内で完結する活動です。}

\quizq{3}{AIの「自社開発」と「外部サービス利用」の比較として適切なものはどれでしょうか?}%
  {\correct{外部サービスは早く導入できるが、差別化やデータの扱いに制約がある場合がある}}%
  {両者に違いはない}%
  {自社開発は常に安い}%
  {外部サービスは常に劣る}%
  {正解は (a)。外部サービスは導入が早く手軽ですが、他社も使えるため差別化しにくく、データの扱いにも制約がある場合があります。自社開発は差別化しやすい反面、コストと人材が必要です。(c)(d) のように一方が常に優れているわけではなく、業務の重要度と自社の強みで使い分けます。}

\quizq{4}{AI投資を評価する際に考慮すべきこととして適切なものはどれでしょうか?}%
  {流行しているかどうかだけ}%
  {競合が導入したかどうかだけ}%
  {\correct{導入・運用コストと期待効果、リスク(データ・法規制など)}}%
  {初期費用だけ}%
  {正解は (c)。AI投資は、導入・運用コストと期待効果に加えて、データの確保や法規制などのリスクも含めて総合的に評価します。(a)(b) 流行や競合の動向だけで決めると目的が曖昧になり、(d) 初期費用だけでは運用コストを見落とします。}

\quizq{5}{AIによるビジネスモデル変化の例として適切なものはどれでしょうか?}%
  {価格を固定する}%
  {全員に同じ商品を郵送する}%
  {\correct{利用データをもとに、個々の顧客に合わせた提案を行うサービス}}%
  {店舗の数を増やす}%
  {正解は (c)。AIは利用データから個々の顧客の好みを推定し、パーソナライズされた提案を可能にします。これは「全員に同じものを売る」モデルから「一人ひとりに合わせる」モデルへの変化です。(a)(b)(d) はAIによるビジネスモデルの変化ではありません。}
